\chapter{Die Notfall-Karte}
\thispagestyle{plain}
Diese Seite ist einzeln gedacht: Drucken Sie sie ein zweites Mal und legen Sie sie sichtbar zum Rechner.
\small

\begin{tikzpicture}
\node[draw=warnrot,line width=2.2pt,rounded corners=4pt,inner sep=12pt,text width=\dimexpr\linewidth-24pt-4.4pt\relax,align=flush left]{%
{\sffamily\bfseries\Large\color{warnrot} NOTFALL-KARTE NOMAD}\par\vspace{6pt}
\textbf{1. Adresse:} Im Browser (Firefox) \texttt{http://localhost:8080} eintippen.\par
Von einem anderen Gerät im selben WLAN: \texttt{http://} \feld[3.8cm] \texttt{:8080}\par\vspace{6pt}
\textbf{2. Der Rechner fährt nicht hoch:} USB-SSD fest eingesteckt? Anderer blauer USB-Anschluss. Startmenü-Taste: \taste{\feld[1.1cm]}\par\vspace{6pt}
\textbf{3. Seite lädt nicht, obwohl der Desktop da ist:} 2 Minuten warten. Dann Terminal (\taste{Strg}+\taste{Alt}+\taste{T}) und:\par
\befehl{sudo bash /opt/project-nomad/start\_nomad.sh}\par
Passwort eingeben (Tipp unsichtbar), \taste{Enter}.\par\vspace{6pt}
\textbf{4. Alles prüfen:} \befehl{sudo docker ps}
Es sollten Zeilen mit \texttt{nomad\_} erscheinen.\par\vspace{6pt}
\textbf{5. Ausschalten:} Oben rechts auf das Ein/Aus-Symbol, \menue{Ausschalten}. Nie die USB-SSD abziehen, solange der Rechner läuft.\par\vspace{6pt}
\textbf{6. Strom sparen im Notfall:} Bildschirm dunkel, WLAN/Bluetooth aus, nur bei Bedarf einschalten (Kapitel~8).\par\vspace{6pt}
\textbf{7. Wissen suchen:} \menue{Wissensbibliothek} (\texttt{http://localhost:8090}), Suchfeld links oben. Texte sind überwiegend englisch.\par\vspace{6pt}
\textbf{8. Passwort:} \feld[6cm] \quad (Ort der Notiz: \feld[3.8cm])\par\vspace{6pt}
\textbf{9. Hilfe:} \feld[10cm]};
\end{tikzpicture}

\vspace{0.5em}
Software-Version 1.35.2, Stand des Heftes Oktober 2026.
